% =============================================================================
% PAPER_FLUX_manuscript_section6_patch.tex
%
% Paste-ready patch for the manuscript
%   "Two-Sided Spectral Control of Wilson One-Plaquette Class Gaps and an
%    Exactly Verified Strong-Coupling Bridge to the SU(3) Glueball Channels"
% implementing master-document items SS6.9 (Theorem 6.2 correction) and SS6.10
% (Theorem 6.3' and the O(y^3) extension), plus the v2.2 ratio corrections and
% the quantum-number relabeling.
%
% Every constant below is drawn from the certified set:
%   ENGINE_FLUX_su3_domino_d3.py (251 gates) / RUN_TROM_d3_results.json   -- t_+, leakage, d_3, all
%       order-2 and order-3 constants;
%   ENGINE_FLUX_glueball_band_certificate_v2.py (36 gates)       -- band structure, flat
%       band, curvature, quantum numbers, corrected C-even block;
%   ENGINE_TROM_tromino_contract_independent_check.py (19 gates),
%   ENGINE_TROM_tromino_candidate_closed_form_check.py           -- O(y^3) tromino sector;
%   ENGINE_FLUX_su3_moments_ext.py (27 gates)                    -- Haar moment engine;
%   ENGINE_FLUX_master_v2_regression_certificate.py (40 gates)   -- assembly arithmetic.
%
% This file compiles standalone (pdflatex) so the blocks can be proofread in
% rendered form; each block is delimited by PATCH comments stating exactly
% where it goes in the manuscript source.
% =============================================================================
\documentclass[11pt]{article}
\usepackage{amsmath,amssymb,amsthm}
\usepackage[margin=1.1in]{geometry}
\newtheorem{theorem}{Theorem}[section]
\newtheorem{lemma}[theorem]{Lemma}
\newtheorem{remark}[theorem]{Remark}
\setcounter{section}{6}
\newcommand{\Tr}{\operatorname{Tr}}
\title{Section 6 patch: corrected leakage constants, the exact band structure,
and the $O(y^3)$ extension}
\author{Patch note compiled against the June 11 certificate chain}
\date{June 11, 2026}
\begin{document}
\maketitle

\noindent\textbf{Summary of changes.}
(i) The C-even hopping of Theorem~6.2 omits a vacuum-mediated second-order
route present for every C-even pair; the corrected per-neighbor hop is
$t_+=-\tfrac{481}{612}+\tfrac34=-\tfrac{11}{306}$, and the $k=0$ coefficient
$-\tfrac{9397}{1020}$ becomes $-\tfrac{217}{1020}$. The manuscript's own
Section~7 domino data adjudicate the correction. (ii) The open signed-band
problem of Theorem~6.3 is solved: the lowest C-odd branch is \emph{exactly
flat} at $O(y^2)$, with coefficient $\tfrac{11}{306}$, and remains exactly
flat at $O(y^3)$ with cubic coefficient $-\tfrac{109151}{249696}$.
(iii) Remark~6.4's per-bond ratio $5/481$ does not survive the correction
(the corrected ratio is $5/22$); the immobility statement survives in the
stronger flat-band form. (iv) The C-odd channel is $T_1^{+-}$
(axial-vector-like), not scalar; the C-even rest-frame content is
$A_1^{++}\oplus E^{++}$ with the new $E^{++}$ coefficient $\tfrac{223}{1020}$.

% =============================================================================
% PATCH 1 of 7  [SS6.9]
% LOCATION: immediately after Lemma 6.1 (after eq. (19) and the channel-energy
% sentence), before Theorem 6.2.
% ACTION: insert the following lemma verbatim.
% =============================================================================
\medskip\hrule\medskip
\noindent\textbf{PATCH 1 --- insert after Lemma 6.1:}

\begin{lemma}[Vacuum-mediated route; $6.1'$]\label{lem:vacroute}
For C-even single-flux states the second-order effective Hamiltonian contains
the vacuum-intermediate route
\begin{equation}
\langle e_r|\,W\,R\,W\,|e_i\rangle \;\supset\;
\frac{\langle e_r|W|0\rangle\,\langle 0|W|e_i\rangle}{\,8/3-0\,}
=\bigl(\sqrt2\bigr)\bigl(\sqrt2\bigr)\tfrac38=+\tfrac34
\tag{$19'$}
\end{equation}
for \emph{every} ordered pair $(i,r)$, adjacent or distant, with the same
normalization as the operator identity displayed in Theorem~6.2
($\langle G_{p'}|W_p=(\chi_p+\bar\chi_p)(\chi_{p'}+\bar\chi_{p'})\langle
0|/\sqrt2$). It is absent in the C-odd sector, where
$\langle o|W|0\rangle=0$ by parity. For a \emph{distant} pair the route
exactly cancels the free two-flux channel sum
$2/(8/3-16/3)=-\tfrac34$, recovering the vanishing of distant hops; for
shared-link neighbors the two-flux sum is the fused-channel value
$-\tfrac{481}{612}$ of Theorem~6.2 and the cancellation is partial.
\end{lemma}

\begin{proof}
The amplitude $\langle e_r|W|0\rangle=\sqrt2$ per flux state follows from the
displayed identity; the intermediate vacuum sits at energy $0$ below the
single-flux manifold at $8/3$, giving the denominator factor $3/8$. For a
distant pair the only other intermediates are the free two-flux states at
$16/3$, whose summed squared amplitude is $2$; the two routes cancel. The
mechanism, both cancellations, and the adjacent-pair value are
machine-certified (su3\_domino\_d3.py, 251 exact-rational gates), and the
Section~7 exact diagonalization adjudicates independently: the C-even domino
levels $\{1769/3060,\,13/20\}$ are
$\text{diag}\pm|t_+|$ with $\text{diag}=1879/3060$, forcing
$|t_+|=110/3060=11/306$, not $481/612$.
\end{proof}

% =============================================================================
% PATCH 2 of 7  [SS6.9]
% LOCATION: Theorem 6.2, equations (20)-(22).
% ACTION: replace the statement of Theorem 6.2 in full by the following.
% The channel split, the orientation independence, the operator identity, the
% vacuum subtraction, and the neighbor count are unchanged; only the hopping
% and the assembled constants change.
% =============================================================================
\medskip\hrule\medskip
\noindent\textbf{PATCH 2 --- replace Theorem 6.2 (eqs.\ (20)--(22)) with:}

\begin{theorem}[C-even leakage; corrected]\label{thm:62corr}
Per shared-link neighbor, in the C-even sector:
\begin{equation}
\text{self-energy}=-\frac{481}{612}
\qquad\Bigl(\text{channel split: }
1:-\tfrac1{12},\;8:-\tfrac{16}{51},\;\bar3:-\tfrac16,\;6:-\tfrac29\Bigr),
\tag{20}
\end{equation}
independent of the relative link orientation, with vacuum subtraction $+3/4$
per neighbor on the diagonal. The hopping is the fused-channel sum
\emph{completed by the vacuum-mediated route of Lemma~$6.1'$}:
\begin{equation}
t_+=-\frac{481}{612}+\frac34=-\frac{11}{306}.
\tag{$20'$}
\end{equation}
The operator identity
$\langle G_{p'}|W_p=\langle G_p|W_{p'}$ equates the \emph{channel sums} of
hopping and self-energy; the full hop differs from the self-energy by the
$|0\rangle$ route, which contributes off-diagonally but is removed from the
diagonal by the vacuum subtraction. Assembling at $k=0$ in the uniform
$A_1^{++}$ projection with neighbor count $4(2d-3)$ in $d$ spatial
dimensions,
\begin{equation}
m_+(y)=4\,\Delta_+\!\Bigl(\tfrac{3y}{2}\Bigr)+y^2 L(d)+O(y^3),
\qquad
L(d)=4(2d-3)\Bigl(-\frac{11}{153}\Bigr),
\tag{21}
\end{equation}
and in particular, for $d=3$,
\begin{equation}
m_+(y)=\frac83-y-\frac{217}{1020}\,y^2+O(y^3),
\qquad \frac{217}{1020}=0.2127\ldots
\tag{22}
\end{equation}
\end{theorem}

\noindent\emph{[Optional one-line addition to Assumption 6.5, after ``the
vacuum path contributing $+3/4$ exactly once'':] ``; its off-diagonal
counterpart is the route of Lemma~$6.1'$.''}

% =============================================================================
% PATCH 3 of 7  [SS6.10]
% LOCATION: immediately after Theorem 6.3 (after eq. (23) and the paragraph
% ending "...fixes the coefficient to [-3/102, 17/102]").
% ACTION: insert the following theorem. It SOLVES the band problem Theorem 6.3
% leaves open; the final sentence of Theorem 6.3's trailing paragraph
% ("...is left open; the bound above already fixes...") should be amended to
% "...is solved in Theorem 6.3'; the bound above already fixes...".
% =============================================================================
\medskip\hrule\medskip
\noindent\textbf{PATCH 3 --- insert after Theorem 6.3:}

\begin{theorem}[Exact $O(y^2)$ band structure; flat C-odd band; $6.3'$]
\label{thm:63prime}
Plaquettes of the cubic lattice carry three orientations, so the
one-excitation effective Hamiltonian at $O(y^2)$ is a $3\times3$ Bloch
problem at each $k$. Let $A(k)$ and $S(k)$ be the unsigned and signed Bloch
matrices of the shared-link adjacency (12-regular; signs from the relative
orientation $s=\pm1$ of Theorem~6.3), and let $N(k),\widetilde N(k)$ be the
Bloch incidence matrices of the oriented boundary chains, with
$u_j=1-e^{ik_j}$, $v_j=1+e^{ik_j}$. Then
\begin{equation}
A(k)+4I=N(k)N(k)^\dagger,\quad \det N(k)=-2\,v_1v_2v_3;
\qquad
S(k)+4I=\widetilde N(k)\widetilde N(k)^\dagger,\quad
\det \widetilde N(k)\equiv 0,
\tag{$23'$}
\end{equation}
so $\lambda(k)\in[-4,12]$, $\mu(k)\in[-4,8]$, and the C-odd spectrum contains
$\mu=-4$ at every $k$, with kernel vector
$w(k)=(\bar u_3,-\bar u_2,\bar u_1)$: the lowest C-odd branch is
\emph{exactly flat},
\begin{equation}
m_-(k,y)=\frac83+y+\frac{11}{306}\,y^2+O(y^3)
\qquad\text{for every }k,
\tag{$23''$}
\end{equation}
replacing the interval $[-\tfrac3{102},\tfrac{17}{102}]$ of Theorem~6.3 by
the exact value $\tfrac{11}{306}=\tfrac{7}{102}-\tfrac{20}{612}$. The flat
band is spanned by compactly localized elementary-cube $2$-chains
($\partial_2\partial_3=0$); on an $L^3$ torus the $\mu=-4$ eigenspace has
dimension $L^3+2$. The dispersive C-odd partner branches span
$\mu\in[-4,8]$, i.e.\ coefficients $[\tfrac{11}{306},\tfrac{41}{306}]$, with
$\mu=-4+|k|^2+O(k^3)$ (doubly degenerate) near rest. In the C-even sector,
with the corrected hop of Theorem~6.2,
\begin{equation}
E_+(k)=\frac83-y+y^2\Bigl[\frac{223}{1020}-\frac{11}{306}\,\lambda(k)\Bigr],
\qquad \lambda(k)\in[-4,12],
\tag{$23'''$}
\end{equation}
with band bottom $-\tfrac{217}{1020}$ at $k=0$ ($\lambda=12$,
$\lambda(k)=12-\tfrac43|k|^2+O(k^4)$ isotropically, hence curvature
$+\tfrac{22}{459}|k|^2y^2$ and effective mass $m^*=459/(44\,y^2)$), band top
$\tfrac{1109}{3060}$ ($\lambda=-4$), and bandwidth
$16|t_+|=\tfrac{88}{153}$. At rest the C-even triplet decomposes as
$A_1^{++}\oplus E^{++}$: the $A_1^{++}$ ($0^{++}$) state is the band bottom,
and the $E^{++}$ doublet ($2^{++}$-like) sits at $\lambda=0$ with exact,
hop-independent coefficient $\tfrac{223}{1020}$. The C-odd rest triplet is
$T_1^{+-}$ ($1^{+-}$-like): \emph{the C-odd one-plaquette sector contains no
scalar at rest.}
\end{theorem}

% =============================================================================
% PATCH 4 of 7  [SS6.10]
% LOCATION: immediately after Theorem 6.3' (Patch 3).
% ACTION: insert the following theorem (the O(y^3) extension).
% =============================================================================
\medskip\hrule\medskip
\noindent\textbf{PATCH 4 --- insert after Theorem $6.3'$:}

\begin{theorem}[Third order: tromino vanishing and the rigid flat band;
$6.3''$]\label{thm:63dprime}
At $O(y^3)$ the only new lattice geometry is the tromino (three plaquettes
pairwise sharing links). \emph{Every $O(y^3)$ tromino weight vanishes}
(bare-link lemma): each candidate three-plaquette process leaves at least one
boundary link bare, and the corresponding Haar integral vanishes; trominoes
first activate at $O(y^4)$. Consequently the third-order effective
Hamiltonian retains the second-order signed-adjacency structure, with C-odd
hopping odd in $s$ ($T_3(s)=-T_3(-s)$), and the flat band of
Theorem~$6.3'$ is rigid at third order:
\begin{equation}
m_-(k,y)=\frac83+y+\frac{11}{306}\,y^2-\frac{109151}{249696}\,y^3+O(y^4)
\qquad\text{for every }k,
\tag{$23''''$}
\end{equation}
with $d_3=\tfrac{7}{32}+12\,\mathrm{leak}_3-4\,b_3=-\tfrac{109151}{249696}
\approx-0.437135$, where $b_3=\tfrac{1975}{124848}$ is the third-order C-odd
hop and $\mathrm{leak}_3=-\tfrac{12331}{249696}$ the per-neighbor diagonal
leakage; the C-odd dispersive top has cubic coefficient
$-\tfrac{61751}{249696}$. In the C-even sector at rest,
\begin{equation}
m_+(0,y)=\frac83-y-\frac{217}{1020}\,y^2-\frac{54049}{520200}\,y^3+O(y^4),
\tag{$23'''''$}
\end{equation}
and at the band top ($\lambda=-4$) the cubic coefficient is
$\tfrac{471353}{1560600}$. The diagonal and hopping corrections obey the
exact identity $\mathrm{leak}^{\mathrm e}_n=T^{\mathrm e}_n$ at $n=2,3$ (two
distinct $+3/4$ mechanisms of equal value), so the $A_1^{++}$ correction at
$k=0$ is (within-tower) $+\,24\,T^{\mathrm e}_n$ per order.
\end{theorem}

% =============================================================================
% PATCH 5 of 7  [SS6.9 + v2.2 ratios]
% LOCATION: Remark 6.4.
% ACTION: replace Remark 6.4 in full by the following.
% =============================================================================
\medskip\hrule\medskip
\noindent\textbf{PATCH 5 --- replace Remark 6.4 with:}

\begin{remark}[An exactly immobile excitation; corrected]\label{rem:64corr}
With the corrected hop the per-bond ratio is
$|t_-|/|t_+|=(5/612)/(11/306)=5/22$, so the ratio no longer carries an
immobility claim. The claim holds in a stronger, exact form
(Theorems~$6.3'$--$6.3''$): the lowest C-odd branch is \emph{exactly
dispersionless} through $O(y^3)$ --- bandwidth $0$, infinite effective mass
at these orders --- while the full C-odd three-branch manifold has width
$\tfrac{5}{51}$ in $y^2$ units against the C-even bandwidth
$\tfrac{88}{153}$ (manifold-width ratio
$(5/51)/(88/153)=15/88$). The flatness is structural
($\partial_2\partial_3=0$ on the signed incidence complex), not a
small-parameter accident. This is, to our knowledge, a new structural
observation about the strong-coupling $T_1^{+-}$ sector.
\end{remark}

% =============================================================================
% PATCH 6 of 7  [global relabeling]
% ACTION: global edits (search-and-replace with judgment):
%   (a) "C-odd scalar channel" / "strong-coupling scalar sector" (where the
%       C-odd channel is meant, e.g. old Remark 6.4) -> "C-odd T_1^{+-}
%       channel" / "strong-coupling T_1^{+-} sector". The assembled C-odd
%       mass is the mass of the degenerate T_1^{+-} triplet, consistent with
%       the lightest C-odd glueball being 1^{+-} in strong coupling.
%   (b) Abstract and Section 1: "C-even hopping -481/612" ->
%       "C-even hopping -11/306 (the fused-channel sum -481/612 completed by
%       a universal vacuum-mediated route +3/4)"; add after the leakage
%       sentence: "the C-odd band is exactly flat through O(y^3)".
%   (c) Everywhere "-9397/1020" -> "-217/1020"; "L(d) = 4(2d-3)(-503/612)"
%       -> "L(d) = 4(2d-3)(-11/153)"; curvature "481/459" -> "22/459";
%       bandwidth "1924/153" -> "88/153"; ratio "5/481 ~ 1/96" -> "5/22".
%   (d) Section 7's table and prose are UNCHANGED (they were correct and are
%       the adjudicator); add one sentence at the end of Section 7: "The
%       C-even level pair {1769/3060, 13/20} is diag +- |t_+| with
%       diag = 1879/3060, which forces |t_+| = 11/306 and adjudicates the
%       correction of Theorem 6.2."
% =============================================================================

% =============================================================================
% PATCH 7 of 7  [Appendix A additions]
% LOCATION: Appendix A certificate list.
% ACTION: append the following items.
% =============================================================================
\medskip\hrule\medskip
\noindent\textbf{PATCH 7 --- append to the Appendix A certificate list:}

\begin{itemize}
\item \texttt{su3\_moments\_ext.py} (27 gates) --- extended exact SU(3) Haar
moment engine: Weingarten $p=q\le3$, charge-$\pm3$ $\varepsilon$-projector
blocks, independent Weyl-torus oracle.
\item \texttt{su3\_domino\_d3.py} (251 gates) --- abstract-domino des
Cloizeaux calculus to $O(y^3)$: the vacuum-route lemma, $t_+=-11/306$, the
Section~7 levels as exact gates, $b_3$, $\mathrm{leak}_3$,
$d_3=-109151/249696$, and the C-even cubic constants.
\item \texttt{glueball\_band\_certificate\_v2.py} (36 gates) --- the
$O(y^2)$ band structure of Theorem~$6.3'$: incidence factorizations,
$\det\widetilde N\equiv0$, cube $2$-chains and the $L^3+2$ multiplicity,
corrected C-even block, $O_h$ quantum numbers; the as-written
$-9397/1020$ is retained as a provenance gate.
\item \texttt{tromino\_contract\_independent\_check.py} (19 gates) and
\texttt{tromino\_candidate\_closed\_form\_check.py} --- the $O(y^3)$
tromino-vanishing sector of Theorem~$6.3''$.
\end{itemize}

\bigskip\hrule\bigskip
\section*{Constant provenance table}
\noindent Every constant in the patches, with its certificate:
\begin{center}
\begin{tabular}{lll}
constant & value & certificate \\\hline
$t_+$ & $-11/306$ & su3\_domino\_d3.py; \S7 adjudication \\
$A_1^{++}$ at $k=0$, $O(y^2)$ & $-217/1020$ & d3\_results.json; band v2 \\
band top, $O(y^2)$ & $1109/3060$ & d3\_results.json; band v2 \\
$E^{++}$ at $k=0$ & $223/1020$ & band v1 and v2 ($t$-independent) \\
curvature & $22/459$ & band v2; Bloch derivation (master \S3.15) \\
bandwidth & $88/153$ & band v2 \\
flat band, $O(y^2)$ & $11/306$ & band v1/v2; su3\_domino\_d3.py \\
$b_3$ & $1975/124848$ & su3\_domino\_d3.py \\
$\mathrm{leak}_3$ (C-odd) & $-12331/249696$ & su3\_domino\_d3.py \\
$d_3$ & $-109151/249696$ & su3\_domino\_d3.py \\
C-odd dispersive top, $O(y^3)$ & $-61751/249696$ & su3\_domino\_d3.py \\
$m_+(0)$ cubic & $-54049/520200$ & su3\_domino\_d3.py \\
C-even band-top cubic & $471353/1560600$ & su3\_domino\_d3.py \\
per-bond ratio & $5/22$ & band v2 \\
manifold-width ratio & $15/88$ & band v2 \\
\end{tabular}
\end{center}
\end{document}
